\begin{ouverture}[piles]{19}{Circuits et lois électriques}{Ondes et signaux}{Observer · Modéliser · Expérimenter · Expliquer}
\objectif{Mesurer intensité et tension avec les branchements adaptés.}
\objectif{Appliquer les lois des nœuds et des mailles.}
\objectif{Tracer une caractéristique et exploiter la loi d’Ohm.}
\prerequis{Proportionnalité, conversions et lecture de graphiques.}
\end{ouverture}
\essentiel
\essSection{1}{Deux grandeurs distinctes}
L'intensité mesure le débit de charge à travers une section du circuit ; la tension compare les potentiels de deux points. Un ampèremètre se branche en série dans la branche étudiée. Un voltmètre se branche en dérivation entre les deux points. Un courant conventionnel sort de la borne positive du générateur dans le circuit extérieur.
\schema[corrige]{mesures-circuit}
On note $U_{AB}=V_A-V_B$ : inverser les bornes change le signe. Pour une mesure positive, les entrées et les polarités du multimètre doivent correspondre au sens choisi.
\essSection{2}{Nœuds et mailles}
\begin{loi}[Loi des nœuds]La somme des intensités entrant dans un nœud égale celle des intensités qui en sortent. Des signes peuvent être utilisés si les flèches ont été choisies arbitrairement.\end{loi}
\schema[corrige]{noeud}
\begin{loi}[Loi des mailles]La somme algébrique des tensions dans une maille est nulle. Dans un circuit série simple, $U_G=U_1+U_2$ avec les polarités indiquées. Deux dipôles branchés entre les mêmes nœuds ont la même tension.\end{loi}
L'intensité ne se consomme pas dans une résistance. La description doit distinguer la circulation de charge et les transferts d'énergie.
\essSection{3}{Caractéristique d’un dipôle}
Une caractéristique relie une tension à une intensité avec des conventions précisées. Pour un conducteur ohmique dans les conditions étudiées, les points $U=f(I)$ se répartissent près d'une droite passant par l'origine.
\begin{loi}[Loi d'Ohm]En convention récepteur, $U=RI$, avec $U$ en volts, $I$ en ampères et $R$ en ohms. La pente de $U=f(I)$ donne $R$ si l'intensité est en ampères.\end{loi}
\schema[corrige]{ohm}
Le point de fonctionnement correspond à un couple $(I;U)$ compatible avec le dipôle et l'alimentation. Tout dipôle n'est pas ohmique : une courbe qui ne passe pas par l'origine ou qui n'est pas droite ne peut pas être résumée par une résistance constante.
\essSection{4}{Expérimenter et modéliser}
\begin{experience}[Relever une caractéristique]Utiliser une alimentation scolaire de très basse tension continue. Varier doucement la tension, relever U en dérivation et I en série, sans dépasser les caractéristiques des composants. Tracer U en fonction de I, convertir les milliampères en ampères avant de calculer la pente. Couper l'alimentation avant de modifier un branchement ; ne jamais relier un ampèremètre seul aux bornes du générateur.\end{experience}
\exercices
\rubrique{Prendre confiance}
\begin{blocLivre}{ex-19-1}\begin{exercice}[1]{Choisir le branchement}
On veut mesurer la tension d’une résistance et l’intensité qui la traverse.
\schema{mesures-circuit}
\begin{questions}
\item Placer voltmètre et ampèremètre.
\item Quel branchement dangereux faut-il éviter pour l’ampèremètre ?
\end{questions}
\end{exercice}\end{blocLivre}
\begin{blocLivre}{ex-19-2}\begin{exercice}[1]{Convertir une intensité}
Un courant a une intensité de 25,0 mA.
\begin{questions}
\item Exprimer I en ampères.
\item Calculer U pour $R=200\,\Omega$.
\end{questions}
\end{exercice}\end{blocLivre}
\begin{blocLivre}{ex-19-3}\begin{exercice}[1]{Une dérivation}
Un courant de 0,30 A se sépare en deux branches ; la première reçoit 0,12 A.
\schema{noeud}
\begin{questions}
\item Calculer le courant de la seconde branche.
\item Citer la loi utilisée.
\end{questions}
\end{exercice}\end{blocLivre}
\begin{blocLivre}{ex-19-4}\begin{exercice}[1]{Une maille simple}
Un générateur de 9,0 V alimente deux dipôles en série. Le premier a une tension de 3,5 V.
\begin{questions}
\item Déterminer la tension du second.
\item La tension est-elle nécessairement la même sur les deux ?
\end{questions}
\end{exercice}\end{blocLivre}
\rubrique{S’entraîner}
\begin{blocLivre}{ex-19-5}\begin{exercice}[2]{Une caractéristique ohmique}
Lire la caractéristique U=f(I) fournie.
\schema{ohm}
\begin{questions}
\item Déterminer la résistance.
\item Lire U pour $I=20$ mA.
\end{questions}
\end{exercice}\end{blocLivre}
\begin{blocLivre}{ex-19-6}\begin{exercice}[2]{Résistance inconnue}
Une tension de 6,0 V produit une intensité de 15 mA dans un conducteur ohmique.
\begin{questions}
\item Calculer R.
\item Quelle intensité prévoir à 3,0 V si R reste constant ?
\end{questions}
\end{exercice}\end{blocLivre}
\begin{blocLivre}{ex-19-7}\begin{exercice}[2]{Deux branches parallèles}
Deux résistances de 100 et 200 ohms sont branchées chacune sous 6,0 V.
\schema{parallele}
\begin{questions}
\item Calculer les intensités des branches.
\item Calculer l’intensité totale.
\end{questions}
\end{exercice}\end{blocLivre}
\begin{blocLivre}{ex-19-8}\begin{exercice}[2]{Inverser les bornes}
Un voltmètre affiche $U_{AB}=+4,2$ V.
\begin{questions}
\item Quelle valeur donnerait $U_{BA}$ ?
\item Le changement de signe indique-t-il une panne ?
\end{questions}
\end{exercice}\end{blocLivre}
\rubrique{Raisonner et approfondir}
\begin{blocLivre}{ex-19-9}\begin{exercice}[3]{Deux résistances en série}
Un générateur idéal de 9,0 V alimente en série $R_1=100\,\Omega$ et $R_2=200\,\Omega$.
\schema{serie}
\begin{questions}
\item Exprimer I avec la loi des mailles et la loi d’Ohm.
\item Calculer I, $U_1$ et $U_2$.
\end{questions}
\end{exercice}\end{blocLivre}
\begin{blocLivre}{ex-19-10}\begin{exercice}[3]{Un courant non consommé}
Un ampèremètre mesure 20 mA avant une résistance dans une unique boucle.
\begin{questions}
\item Quelle intensité doit-on mesurer après elle en régime établi ?
\item Pourquoi « la résistance consomme le courant » est-il incorrect ?
\end{questions}
\end{exercice}\end{blocLivre}
\begin{blocLivre}{ex-19-11}\begin{exercice}[3]{Un dipôle non ohmique}
Les couples $(I;U)$ sont (10 mA ; 1,0 V), (20 mA ; 2,8 V) et (30 mA ; 5,5 V).
\begin{questions}
\item Comparer les quotients U/I.
\item Le modèle d’une résistance constante convient-il ?
\end{questions}
\end{exercice}\end{blocLivre}
\begin{blocLivre}{ex-19-12}\begin{exercice}[3]{Une modélisation numérique}
Un ajustement donne $U=198I+0,02$ avec U en V et I en A. Les écarts expérimentaux sont de l’ordre de 0,03 V.
\begin{questions}
\item Quelle résistance approchée suggère la pente ?
\item Peut-on considérer le petit décalage comme nécessairement significatif ?
\end{questions}
\end{exercice}\end{blocLivre}

